% Author contribution and AI-use disclosure. Placement: after the appendix, before the
% bibliography.
%
% Form (the author's decision of 2026-09-14, response-plan D5): a short contributor section --
% roles, verification and review status, responsibility -- in the first person, since it is the
% author's own declaration (ADR-0015's neutral third person governs the body). The chronology,
% the figures with their caveats, and the detailed account of the reviews are in the
% repository's notes/PROCESS.md, which this section cites; that note is the material for any
% longer account. The hub standard's four rules (writing-style-and-ai-use.md section 6) apply:
% never claim less involvement than there was; numbers from chronicler with their date and
% caveats (in the note, not here); do not repeat the trust boundary, which section 1.1 owns;
% anonymity outranks completeness (not in play, the release is not double-blind).
%
% History: the first version (2026-09-12) was on Paper I's model with the chronology and the
% numbers in the paper; the author's review (2026-09-14) put it in the first person and brought
% it up to date; the external reviews of the same day both asked for the chronology and the
% figures to move to supplementary material, and the author agreed. What moved is in
% notes/PROCESS.md verbatim, including the three places where the archive corrected the
% author's recollection. Update the note, not this section, when the figures are re-derived.

\section*{Author contribution and use of AI}

This article was written in close collaboration with generative AI, and the collaboration
went beyond editing assistance at every stage. I write this section in the first person
because it is my own declaration; the article itself is written in the neutral third person.

\paragraph{Roles.} The research programme and the question are mine. The axioms of
\sscite{pauwels1995extended} were the template for my semigroup theory of temporal scale
space \sscite{fagerstrom2005temporal}, and after the time-causal article had been drafted I
asked what the hemigroup weakening would give on the line, naming that paper as the point of
departure. Every decision of scope, method and policy is mine: the method of drafting, one
section per exchange in dialogue; the notation; the reference policy toward the companion
article; the split of the material into three modules and the decision to release this one
first; the cut of this version at the characterization theorem; the rule that the proof of
record follows the machine-checked route; and the voice of the text. I reviewed the article
section by section and set the scope of every claim.

The prose and the code were written by AI agents throughout, in dialogue with me; I worked
almost entirely from the conversation rather than from the editor. The Lean~4 development is
theirs essentially in full, and none of the proof terms are mine. So is mathematical content:
the results this article rests on were identified by an agent on reading the source, before
any plan existed, and during the work the agents found and repaired errors in the mathematics
of record that no human had seen. The technical decisions of the formalization were, for the
most part, the agents' recommendations approved by me rather than my own drafts; I read each
with its consequences and followed all of them.

\paragraph{Verification and review.} What is machine-checked, and the axioms on which the
verified development rests, are set out in \S\ref{sec:machine-checked}. Review by agents that
had no part in the writing was a working method throughout: the statement skeleton before any
proof was attempted, the finished development against the article in a fidelity review, and
the sections blind against the writing standard. Before the final revision the manuscript was
read by AI systems of two other vendors, in two rounds: a referee-style review returning a
verdict of major revision and a presentation review, then, on the revised text, a second
presentation review and a referee report the reviewing system left unfinished. Every
finding was verified against the text and the formal development before it was acted on,
and the errors the first referee report found, in the one construction that is not
machine-checked and in claims stated beyond their results, are corrected in this version.
The reviews are archived verbatim in the repository beside the plan that records the
disposition of each finding. Beyond me, no human reader has
yet reviewed the article.

\paragraph{The record.} The sessions were archived as the work went on, so the division of
labour can be stated rather than estimated, and it is stated, with its dates and its caveats,
in the repository's process note (\texttt{notes/PROCESS.md}), together with the chronology of
the work and the places where the archive corrected my own recollection of it. The short
form: the sustained work took eight days, the agents' words outnumber mine by about
twenty-two to one, which is the ratio of the companion article, and the difference from that
article is that the formalization ran as parallel agents with me approving their
recommendations. Those figures are lower bounds. I would not, unaided and as a spare-time
activity, have written this article.

\paragraph{Responsibility.} I directed the work throughout, have satisfied myself of the
mathematical content, and have checked every citation. Responsibility for this article,
including for its errors, is mine alone; the validity of the machine-checked results depends
on the Lean kernel and not on how, or by whom, the proof terms were produced.
